\scanpage{14}
% Pencil foliation: 12, upper right; 13, lower left.
d'où un $G$ schéma en groupes sur $K$. Soit
$I_G$ le ``schéma des invariants'', isomorphe si $G$ connexe
à $T/W$ ($T$ tore maximal de $G$, $W$ le groupe de Weyl).
Alors on trouve un élément $f_s$ de $I_G(\mathbb{Q})$.
(NB. $I_G$ est en fait défini sur $\mathbb{Q}$).
On trouve
ainsi une famille d'éléments $(f_s)_{s\in S^0}$, où $S^0$
désigne l'ens des pts de $S$ tels que $k(s)$ soit
fini : famille des frobenius. Compatibilité avec
les homs
\diagram{\begin{tikzcd}
G\deleted{\unsure{(K)}}\arrow[r]&\pi_1(S,\unsure{\xi|\zeta})
\end{tikzcd}}
et les homs
\diagram{\begin{tikzcd}
\pi_1(S,\unsure{\xi|\zeta})\arrow[r]&G(\mathbb{Q}_\ell).
\end{tikzcd}}
